% GENERATED FILE. Edit canonical lessons, then run npm run generate:books.
% canonical-source-hash: fnv1a64:c9d2f8f0cca01504
% canonical-lessons: GE-C190-vergleichen, GE-C190-verletzen, GE-C190-verpassen, GE-C190-verschieben, GE-C190-verwenden

\chapter{To compare, To injure, To miss, To postpone, To use}
\label{ch:ge-to-compare-to-injure-to-miss-to-postpone-to-use}
\hlchaptermodality{190}

\begin{chapteropening}
\textbf{By the end of this chapter:} \emph{I can say to compare, to injure, to miss, to postpone and to use.}

Complete the last lesson of chapter 190: I can say to compare, to injure, to miss, to postpone and to use.
\end{chapteropening}

\section[vergleichen]{vergleichen — to compare}

\label{lesson:GE-C190-vergleichen}



\begin{quote}
Before the new one: say the German for to forbid, then the German for to connect.
\end{quote}

\subsection*{You'll want to know: vergleichen}
\textbf{vergleichen} — \textquotedblleft{}to compare\textquotedblright{}.

\begin{cousinweb}
\textbf{vergleichen} is \emph{ver-} + \emph{gleich}, \textquotedblleft{}alike\textquotedblright{}, + \emph{-en} (EtymDB 2.1): to set things side by side to see how alike they are.
\end{cousinweb}

\begin{grammarlens}[title={What you've built}]
One more word for this chapter.
\end{grammarlens}

\subsection*{Your turn}
\begin{itemize}
\raggedright
  \item \emph{Say it:} \emph{vergleichen}
  \item \emph{Say it:} \emph{vergleichen}, once more
  \item \emph{Say it:} say \emph{vergleichen}
  \item \emph{Recall:} say the German for to forbid, then the German for to connect, then say \emph{vergleichen} again
\end{itemize}

\begin{tcolorbox}[breakable,colback=teal!4,colframe=teal!35!black,title={Before you move on}]
What does vergleichen mean? (\textquotedblleft{}To compare\textquotedblright{}.) Say it once more.
\end{tcolorbox}
\section[verletzen]{verletzen — to injure}

\label{lesson:GE-C190-verletzen}



\begin{quote}
Before the new one: say the German for to connect, then the German for to compare.
\end{quote}

\subsection*{You'll want to know: verletzen}
\textbf{verletzen} — \textquotedblleft{}to injure\textquotedblright{}.

\begin{cousinweb}
\textbf{verletzen} $\leftarrow$ Middle High German \emph{verletzen}, \textquotedblleft{}to injure\textquotedblright{} (EtymDB 2.1): \emph{ver-} + an old verb \textquotedblleft{}to hinder, to harm\textquotedblright{}, built on the root of English \emph{late}, \textquotedblleft{}slow\textquotedblright{}.
\end{cousinweb}

\begin{grammarlens}[title={What you've built}]
One more word for this chapter.
\end{grammarlens}

\subsection*{Your turn}
\begin{itemize}
\raggedright
  \item \emph{Say it:} \emph{verletzen}
  \item \emph{Say it:} \emph{verletzen}, once more
  \item \emph{Say it:} say \emph{verletzen}
  \item \emph{Recall:} say the German for to connect, then the German for to compare, then say \emph{verletzen} again
\end{itemize}

\begin{tcolorbox}[breakable,colback=teal!4,colframe=teal!35!black,title={Before you move on}]
What does verletzen mean? (\textquotedblleft{}To injure\textquotedblright{}.) Say it once more.
\end{tcolorbox}
\section[verpassen]{verpassen — to miss (a train, a chance)}

\label{lesson:GE-C190-verpassen}



\begin{quote}
Before the new one: say the German for to compare, then the German for to injure.
\end{quote}

\subsection*{You'll want to know: verpassen}
\textbf{verpassen} — \textquotedblleft{}to miss (a train, a chance)\textquotedblright{}.

\begin{cousinweb}
\textbf{verpassen} is \emph{ver-} + the verb \textquotedblleft{}to fit, to pass\textquotedblright{}, from Dutch \emph{passen} $\leftarrow$ French \emph{passer}, \textquotedblleft{}to pass\textquotedblright{} (EtymDB 2.1): to let something pass you by.
\end{cousinweb}

\begin{grammarlens}[title={What you've built}]
One more word for this chapter.
\end{grammarlens}

\subsection*{Your turn}
\begin{itemize}
\raggedright
  \item \emph{Say it:} \emph{verpassen}
  \item \emph{Say it:} \emph{verpassen}, once more
  \item \emph{Say it:} say \emph{verpassen}
  \item \emph{Recall:} say the German for to compare, then the German for to injure, then say \emph{verpassen} again
\end{itemize}

\begin{tcolorbox}[breakable,colback=teal!4,colframe=teal!35!black,title={Before you move on}]
What does verpassen mean? (\textquotedblleft{}To miss (a train, a chance)\textquotedblright{}.) Say it once more.
\end{tcolorbox}
\section[verschieben]{verschieben — to postpone}

\label{lesson:GE-C190-verschieben}



\begin{quote}
Before the new one: say the German for to injure, then the German for to miss.
\end{quote}

\subsection*{You'll want to know: verschieben}
\textbf{verschieben} — \textquotedblleft{}to postpone\textquotedblright{}.

\begin{cousinweb}
\textbf{verschieben} is \emph{ver-} + \emph{schieben}, \textquotedblleft{}to push\textquotedblright{}, the same verb as English \emph{shove} (EtymDB 2.1): to push something to another place or time.
\end{cousinweb}

\begin{grammarlens}[title={What you've built}]
One more word for this chapter.
\end{grammarlens}

\subsection*{Your turn}
\begin{itemize}
\raggedright
  \item \emph{Say it:} \emph{verschieben}
  \item \emph{Say it:} \emph{verschieben}, once more
  \item \emph{Say it:} say \emph{verschieben}
  \item \emph{Recall:} say the German for to injure, then the German for to miss, then say \emph{verschieben} again
\end{itemize}

\begin{tcolorbox}[breakable,colback=teal!4,colframe=teal!35!black,title={Before you move on}]
What does verschieben mean? (\textquotedblleft{}To postpone\textquotedblright{}.) Say it once more.
\end{tcolorbox}
\section[verwenden]{verwenden — to use}

\label{lesson:GE-C190-verwenden}



\begin{quote}
Before the new one: say the German for to miss, then the German for to postpone.
\end{quote}

\subsection*{You'll want to know: verwenden}
\textbf{verwenden} — \textquotedblleft{}to use\textquotedblright{}.

\begin{cousinweb}
\textbf{verwenden} is \emph{ver-} + the verb \textquotedblleft{}to turn\textquotedblright{}, Old High German \emph{wenten} $\leftarrow$ Proto-Germanic \emph{*wandijaną} (EtymDB 2.1), the same verb as English \emph{wend}: to turn something to a use.
\end{cousinweb}

\begin{grammarlens}[title={What you've built}]
One more word for this chapter.
\end{grammarlens}

\subsection*{Your turn}
\begin{itemize}
\raggedright
  \item \emph{Say it:} \emph{verwenden}
  \item \emph{Say it:} \emph{verwenden}, once more
  \item \emph{Say it:} say \emph{verwenden}
  \item \emph{Recall:} say the German for to miss, then the German for to postpone, then say \emph{verwenden} again
\end{itemize}

\begin{tcolorbox}[breakable,colback=teal!4,colframe=teal!35!black,title={Before you move on}]
What does verwenden mean? (\textquotedblleft{}To use\textquotedblright{}.) Say it once more.
\end{tcolorbox}
